\chapter{Le marché des options de change}\label{ch:m2:the-fx-options-market}

Sur la plupart des marchés d'options, un trader demande le prix d'une option à un strike.
Sur le marché des changes, la question est différente. L'écran d'un courtier
pour le dollar--yen à trois mois affiche trois chiffres~: une volatilité à la monnaie,
un \omterm{def:m2:the-fx-options-market:rr}{risk reversal} à 25 delta et un \omterm{def:m2:the-fx-options-market:rr}{butterfly} à 25 delta. Il n'y a pas de strike sur l'écran, et
les chiffres n'ont pas le même sens d'une paire à l'autre~: le fait que «~25 delta~» soit
un \omterm{def:m2:the-fx-options-market:delta}{delta spot} ou un \omterm{def:m2:the-fx-options-market:delta}{delta forward}, le fait qu'il tienne compte de la prime, quel strike est
«~à la monnaie~», tout cela dépend de conventions qui diffèrent entre le dollar--yen et
l'euro--dollar. Deux traders qui évaluent la même option à partir du même écran avec
des conventions différentes obtiennent des strikes et des prix différents. Les options de change
ont plus que doublé leurs volumes entre les enquêtes des banques centrales de 2022 et
2025, pour atteindre 7~\% de l'ensemble des échanges de change, et chacune de ces transactions commence par
à s'accorder sur ces conventions. Ce chapitre les explique, montre comment un smile se
construit à partir des trois cotations, et suit les flux que les \omterm{def:m2:the-fx-options-market:barrier}{options à barrière} envoient
vers le marché spot.

\section{Delta et conventions à la monnaie}

Le modèle retenu est celui de Black--Scholes à deux taux d'intérêt, le
modèle de Garman--Kohlhagen~: pour une paire BASEQUOTE de spot $S$, le taux de la
devise de cotation $r_d$ et celui de la devise de base $r_f$, le forward est
$F = S e^{(r_d - r_f)T}$ et un call sur la devise de base de strike $K$ vaut
$e^{-r_d T}\bigl(F N(d_1) - K N(d_2)\bigr)$, avec les $d_{1,2}$ habituels. Ce qui change
par rapport au livre~1 est le delta, que le marché des changes cote de quatre façons.

\begin{definition}[Delta spot, delta forward et delta ajusté de la prime]\label{def:m2:the-fx-options-market:delta}
Le \emph{delta spot}\index{delta spot} d'une option de change est le montant de la devise de base
par unité de notionnel, à échanger sur le marché spot pour la couvrir~:
$\phi e^{-r_f T} N(\phi d_1)$, avec $\phi = 1$ pour un call et $-1$ pour un put. Le
\emph{delta forward}\index{delta forward} est le montant à échanger sur le marché à terme,
$\phi N(\phi d_1)$. Le \emph{delta ajusté
de la prime}\index{delta ajusté de la prime} tient compte d'une prime payée dans
la devise de base, qui est elle-même une position dans cette devise~:
$\phi e^{-r_f T} (K/F) N(\phi d_2)$ (spot) ou $\phi (K/F) N(\phi d_2)$ (forward).
\end{definition}

L'ajustement compte. Un call en euros dont le delta vaut 60~\% et dont la prime,
73\,669 euros par million, payée en euros, se couvre en n'achetant que 526\,331
euros~: son \omterm{def:m2:the-fx-options-market:delta}{delta ajusté de la prime} vaut 52.63~\%. Le delta qu'utilise une paire est une
convention, documentée dans la littérature sur la construction du smile de change~: les paires
constituées de devises de pays riches utilisent des \omterm{def:m2:the-fx-options-market:delta}{deltas spot} jusqu'à un an et des \omterm{def:m2:the-fx-options-market:delta}{deltas forward} au-delà~;
les paires avec une devise de pays émergent utilisent des \omterm{def:m2:the-fx-options-market:delta}{deltas forward}~; et le delta est
ajusté de la prime lorsque celle-ci est payée dans la devise de base, comme pour le dollar--yen,
le dollar--franc suisse et l'euro--yen, mais pas pour l'euro--dollar ni la livre sterling--dollar,
où elle est payée en dollars, la devise de cotation.

\begin{definition}[Straddle delta neutre]\label{def:m2:the-fx-options-market:dns}
Le \emph{straddle delta neutre}\index{straddle delta neutre} est le call et le
put au strike où leurs deltas s'annulent~; ce strike est la valeur par défaut de la
«~à la monnaie~» pour les options de change à court terme~: $K = F e^{\sigma^2 T/2}$ avec des
deltas classiques et $K = F e^{-\sigma^2 T/2}$ avec des \omterm{def:m2:the-fx-options-market:delta}{deltas ajustés de la prime}.
\end{definition}

\begin{example}[Strikes à trois mois]\label{ex:m2:the-fx-options-market:strikes}
Avec EURUSD à 1.1464, des taux du dollar et de l'euro de 3.68~\% et 2.00~\% (simples, pour
trois mois) et une volatilité de 7~\%, le forward vaut 1.15119 et le
strike du \omterm{def:m2:the-fx-options-market:dns}{straddle delta neutre} 1.15190. Pour USDJPY à 156.87, avec un taux du yen de
0.977~\% et une volatilité de 9.5~\%, le forward vaut 155.8196~; le strike du straddle est
155.9955 avec des deltas classiques et 155.6439 avec des \omterm{def:m2:the-fx-options-market:delta}{deltas ajustés de la prime}, un écart de
35 pips de 0.01 yen entre deux définitions d'une même «~à la monnaie~».
\end{example}

\section{Risk reversals, butterflies et smile}

\begin{definition}[Risk reversal, strangle et butterfly]\label{def:m2:the-fx-options-market:rr}
Le \emph{risk reversal}\index{risk reversal} à 25 delta est la différence entre
la volatilité implicite du call à 25 delta et celle du put à 25 delta,
$\sigma_{25C} - \sigma_{25P}$~; sous forme de transaction, c'est l'achat d'un call et la vente d'un put (ou
l'inverse). Le \emph{strangle}\index{strangle} à 25 delta est un achat de call à 25 delta
et d'un put à 25 delta~; sa volatilité, rapportée à la volatilité à la monnaie, est
le \emph{butterfly}\index{butterfly} à 25 delta.
\end{definition}

Le \omterm{def:m2:the-fx-options-market:rr}{risk reversal} mesure la pente du smile~: il est négatif lorsque les puts de la devise de base
sont plus chers que ses calls, comme lorsqu'une baisse du dollar face au yen est
redoutée plus qu'une hausse. Le \omterm{def:m2:the-fx-options-market:rr}{butterfly} mesure sa courbure~: ce que coûtent les ailes
par rapport au centre. Un courtier qui affirme qu'un \omterm{def:m2:the-fx-options-market:rr}{risk reversal} vaut «~1.2, puts au-dessus~»
signifie que la volatilité du put dépasse celle du call de 1.2 points
(\cref{fig:m2:the-fx-options-market:quotes}). La façon la plus simple
de transformer les trois cotations en deux volatilités d'ailes, souvent utilisée comme première
approximation, est la suivante~:
\begin{equation}\label{eq:m2:the-fx-options-market:simplified}
\sigma_{25C} = \sigma_{\mathrm{ATM}} + \mathrm{BF} + \tfrac12\mathrm{RR}, \qquad
\sigma_{25P} = \sigma_{\mathrm{ATM}} + \mathrm{BF} - \tfrac12\mathrm{RR},
\end{equation}
qui reproduit exactement le \omterm{def:m2:the-fx-options-market:rr}{risk reversal}. Le \omterm{def:m2:the-fx-options-market:rr}{strangle} coté sur le marché est un objet
légèrement différent de la moyenne des deux ailes, et construire un smile
qui le reproduise exactement exige une calibration~; l'approximation retenue est celle que nous
utilisons.

\begin{omfigure}
\begin{tikzpicture}
\begin{axis}[width=9cm,height=4.6cm,
  xlabel={delta},ylabel={volatilité implicite (\%)},
  tick label style={font=\small},label style={font=\small},
  xmin=-35,xmax=35,ymin=8.9,ymax=10.8,grid=major,grid style={gray!25},
  xtick={-25,0,25},xticklabels={put à 25 delta,ATM,call à 25 delta},
  table/col sep=comma]
\addplot[omThm,very thick,mark=*,mark size=2pt] table[x=x,y=vol]{figdata/markets-2/19-the-fx-options-market/quotes.csv};
\draw[omProp,thick,{Stealth}-{Stealth}] (axis cs:29,9.2) -- (axis cs:29,10.4);
\node[omProp,font=\scriptsize,anchor=east,align=right] at (axis cs:28.5,10.35) {RR\\$-1.2$};
\draw[omDef,thick,dashed] (axis cs:-25,9.8) -- (axis cs:25,9.8);
\node[omDef,font=\scriptsize,anchor=south] at (axis cs:6,9.83) {ATM + BF = 9.8};
\end{axis}
\end{tikzpicture}

\omcaption{Les trois cotations de l'\cref{ex:m2:the-fx-options-market:jpy} comme trois points
du smile USDJPY~: l'aile du put à 10.4~\%, à la monnaie 9.5~\%, l'aile du call à 9.2~\%.
Le \omterm{def:m2:the-fx-options-market:rr}{risk reversal} est la différence entre les ailes, le \omterm{def:m2:the-fx-options-market:rr}{butterfly} mesure la
hauteur de leur moyenne au-dessus du centre. Illustratif.}\label{fig:m2:the-fx-options-market:quotes}
\end{omfigure}

\begin{proposition}[De cotations à sourire]\label{prop:m2:the-fx-options-market:smile}
Étant donnés $\sigma_{\mathrm{ATM}}$, RR et BF ainsi qu'une convention de delta, les
volatilités à delta 25 de l'\cref{eq:m2:the-fx-options-market:simplified} et les strikes
$K_{25C}$, $K_{25P}$ auxquels les options dotées de ces volatilités ont des deltas de
$0.25$ et $-0.25$ donnent, avec le strike à la monnaie, trois points du sourire
en strike~; toute interpolation passant par eux, ici un polynôme du second degré en $\log K$,
fournit une volatilité à chaque strike.
\end{proposition}

\begin{proof}
Pour les deltas réguliers, le strike se déduit en forme fermée de
$N(\phi d_1) = \phi\Delta e^{r_f T}$ (spot) ou $\phi\Delta$ (forward)~; pour les
\omterm{def:m2:the-fx-options-market:delta}{deltas ajustés de la prime}, il se calcule numériquement, du côté de la plage de strikes
où le delta est monotone. Trois points distincts déterminent le polynôme du second degré.
\end{proof}

\begin{omfigure}
\begin{tikzpicture}
\begin{axis}[width=11cm,height=5cm,
  xlabel={strike USDJPY},ylabel={volatilité implicite (\%)},
  tick label style={font=\small},label style={font=\small},
  xmin=145,xmax=166,ymin=8.8,ymax=11.8,grid=major,grid style={gray!25},
  legend columns=2,legend style={font=\footnotesize,at={(0.5,-0.3)},anchor=north,draw=none,fill=none,/tikz/every even column/.append style={column sep=8pt}},
  table/col sep=comma]
\addplot[omDef,very thick] table[x=k_reg,y=v_reg]{figdata/markets-2/19-the-fx-options-market/smile_usdjpy.csv};
\addplot[omThm,very thick,dashed] table[x=k_pa,y=v_pa]{figdata/markets-2/19-the-fx-options-market/smile_usdjpy.csv};
\legend{deltas spot réguliers,deltas spot ajustés de la prime}
\end{axis}
\end{tikzpicture}

\omcaption{Un sourire USDJPY à trois mois construit à partir des mêmes trois cotations (ATM 9.5\%,
\omterm{def:m2:the-fx-options-market:rr}{risk reversal} à delta 25 de $-1.2$, \omterm{def:m2:the-fx-options-market:rr}{butterfly} de 0.3) selon deux conventions de delta. La
convention ajustée de la prime, celle du marché pour cette paire, place les strikes des ailes et
du centre plus bas, de sorte qu'une même volatilité correspond à un strike différent.
Illustratif~; données~: tutoriel du chapitre.}\label{fig:m2:the-fx-options-market:smile}
\end{omfigure}

\begin{example}[Deux sourires à partir d'un même écran]\label{ex:m2:the-fx-options-market:jpy}
USDJPY à trois mois~: ATM 9.5\%, \omterm{def:m2:the-fx-options-market:rr}{risk reversal} $-1.2$ (puts en dollars recherchés), \omterm{def:m2:the-fx-options-market:rr}{butterfly}
de 0.3. Par l'\cref{eq:m2:the-fx-options-market:simplified}, le call en dollars à delta 25 vaut
9.2\% et le put 10.4\%. Avec des \omterm{def:m2:the-fx-options-market:delta}{deltas spot} réguliers, leurs strikes sont 160.85
et 150.71~; avec des \omterm{def:m2:the-fx-options-market:delta}{deltas spot} ajustés de la prime, la convention de la paire, 160.68 et
150.52. Un trader qui aurait évalué le put à 150.71 à 10.4\% alors que le marché visait 150.52
commettrait une erreur de prix égale à la pente du sourire sur 19 pips de strike
(\cref{fig:m2:the-fx-options-market:smile}).
\end{example}

\section{Devise de la prime et delta ajusté de la prime}

La devise de la prime modifie à la fois la couverture et la cotation. Une banque
basée en dollars qui vend un call USDJPY et reçoit sa prime en dollars détient déjà une partie des
dollars que sa couverture exige~: elle en achète moins. Le marché cote donc les
deltas de la paire ajustés de la prime, et le \omterm{def:m2:the-fx-options-market:delta}{delta ajusté de la prime} du call présente un maximum~:
profondément dans la monnaie, il croît avec le strike, puis décroît, de sorte que
deux strikes peuvent partager un même delta, et la convention est de retenir celui situé au-dessus du maximum
(\cref{fig:m2:the-fx-options-market:deltas}). Le tutoriel procède ainsi.

\begin{omfigure}
\begin{tikzpicture}
\begin{axis}[width=11cm,height=4.8cm,
  xlabel={strike USDJPY},ylabel={delta spot du call},
  tick label style={font=\small},label style={font=\small},
  xmin=120,xmax=200,ymin=0,ymax=1.05,grid=major,grid style={gray!25},
  legend columns=2,legend style={font=\footnotesize,at={(0.5,-0.3)},anchor=north,draw=none,fill=none,/tikz/every even column/.append style={column sep=8pt}},
  table/col sep=comma]
\addplot[omDef,very thick] table[x=k,y=reg]{figdata/markets-2/19-the-fx-options-market/deltas.csv};
\addplot[omThm,very thick,dashed] table[x=k,y=pa]{figdata/markets-2/19-the-fx-options-market/deltas.csv};
\draw[gray,dashed] (axis cs:141,0) -- (axis cs:141,1.05);
\node[gray,font=\scriptsize,anchor=west,align=left] at (axis cs:166,0.62) {en pointillés~: maximum\\du delta ajusté\\de la prime};
\legend{régulier,ajusté de la prime}
\end{axis}
\end{tikzpicture}

\omcaption{\omterm{def:m2:the-fx-options-market:delta}{Deltas spot} des calls USDJPY à trois mois selon le strike, à une volatilité de 9.5\%. Le
delta régulier décroît régulièrement à partir de un~; le \omterm{def:m2:the-fx-options-market:delta}{delta ajusté de la prime} est plus bas et, profondément
dans la monnaie, croît avec le strike avant de décroître, de sorte qu'un delta situé sous son maximum correspond
à deux strikes. Illustratif~; données~: tutoriel du
chapitre.}\label{fig:m2:the-fx-options-market:deltas}
\end{omfigure} La conséquence pour un trader est pratique~: avant de calculer le moindre
nombre, il faut vérifier le type de delta de la paire, la définition de l'à la monnaie et la
devise de la prime de la cotation.

\section{Banques courtières et flux induits par les barrières}

\begin{definition}[Option à barrière, knock-out, knock-in]\label{def:m2:the-fx-options-market:barrier}
Une \emph{option à barrière}\index{option à barrière} est une option dont l'existence
dépend de la question de savoir si le spot touche un niveau, la barrière, avant l'échéance. Une
\emph{option knock-out}\index{option knock-out} cesse d'exister lorsque la barrière
est touchée~; une \emph{option knock-in}\index{option knock-in} n'entre en existence
qu'à ce moment-là.
\end{definition}

Entre dealers, les options de change vanille se cotent en straddles à la monnaie, \omterm{def:m2:the-fx-options-market:rr}{risk
reversals} et \omterm{def:m2:the-fx-options-market:rr}{strangles} plutôt qu'en strikes isolés, par l'intermédiaire de courtiers et sur des plateformes
électroniques. Les barrières et autres exotiques sont vendues à des entreprises et à des
investisseurs, et couvertes par les dealers en vanilles et en spot. Les barrières séduisent
les acheteurs parce qu'elles coûtent moins cher que les vanilles, et elles comptent pour le marché au comptant
car leurs couvertures changent brusquement à la barrière.

\begin{omfigure}
\begin{tikzpicture}
\begin{axis}[name=left,width=6.4cm,height=4.8cm,
  xlabel={spot EURUSD},ylabel={valeur (pips de notionnel)},
  tick label style={font=\small},label style={font=\small},
  xmin=1.10,xmax=1.17,ymin=-5,ymax=320,grid=major,grid style={gray!25},
  xticklabel style={/pgf/number format/.cd,fixed,precision=2},
  table/col sep=comma]
\addplot[omThm,very thick] table[x=s,y=pips]{figdata/markets-2/19-the-fx-options-market/barrier.csv};
\draw[omProp,dashed] (axis cs:1.12,-5) -- (axis cs:1.12,320);
\node[omProp,font=\scriptsize,anchor=south west] at (axis cs:1.121,250) {barrière};
\end{axis}
\begin{axis}[at={(left.east)},anchor=west,xshift=1.4cm,width=6.4cm,height=4.8cm,
  xlabel={spot EURUSD},ylabel={delta},
  tick label style={font=\small},label style={font=\small},
  xmin=1.10,xmax=1.17,ymin=-0.05,ymax=0.9,grid=major,grid style={gray!25},
  xticklabel style={/pgf/number format/.cd,fixed,precision=2},
  table/col sep=comma]
\addplot[omThm,very thick,const plot] table[x=s,y=delta]{figdata/markets-2/19-the-fx-options-market/barrier.csv};
\draw[omProp,dashed] (axis cs:1.12,-0.05) -- (axis cs:1.12,0.9);
\end{axis}
\end{tikzpicture}

\omcaption{\omterm{def:m2:the-fx-options-market:barrier}{Option knock-out} (down-and-out call) EURUSD à trois mois, prix d'exercice 1.1464, barrière 1.1200,
volatilité 7~\%~: sa valeur (à gauche) et son delta (à droite) en fonction du spot. La valeur tombe
à zéro à la barrière, mais le delta ne le fait pas~: il reste proche de 0.63 tant que
le spot n'a pas touché la barrière, puis tombe à zéro. Illustratif~; données~:
tutoriel du chapitre.}\label{fig:m2:the-fx-options-market:barrier}
\end{omfigure}

\begin{proposition}[La couverture à la barrière]\label{prop:m2:the-fx-options-market:jump}
Un dealer vendeur d'une \omterm{def:m2:the-fx-options-market:barrier}{option knock-out} de notionnel $N$ et couvert en delta détient $\Delta N$
de la devise de base. Lorsque le spot touche la barrière, l'option et son delta
disparaissent, et le dealer doit échanger $\Delta(B^+)\,N$ immédiatement, dans le sens qui
pousse le spot au-delà de la barrière~: pour un call down-and-out, vendre à la
barrière.
\end{proposition}

\begin{proof}
Avant le contact, la couverture vaut $\Delta N$, avec $\Delta \to \Delta(B^+) > 0$ quand le spot
s'approche de $B$ par le haut~; après le contact, l'option a disparu et la couverture doit être nulle.
La couverture d'un call down-and-out est longue en devise de base, donc son dénouement vend à
la barrière, où le spot vient de baisser.
\end{proof}

C'est pourquoi les barrières sont défendues et attaquées. Les détenteurs d'\omterm{def:m2:the-fx-options-market:barrier}{options knock-out}, et les dealers
qui en sont longs, achètent près de la barrière afin qu'elle ne soit pas franchie~; les dealers vendeurs
savent que leurs propres dénouements de couverture aggraveront un mouvement qui la franchit, et ceux qui
savent où se trouvent les grandes barrières peuvent tenter de pousser le spot jusqu'à elles. Les ordres stop peuvent
se placer aux mêmes niveaux, pour la même raison~: ce sont des niveaux que beaucoup de monde surveille.

\begin{dated}{2026-09}{Taille du marché}\label{dat:m2:the-fx-options-market:size}
En avril 2025, les options de change représentaient 7~\% du volume mondial de change, plus du double de leur
volume de 2022, selon le BIS Triennial Survey. L'enquête ne répartit pas
les options vanilles et les \omterm{def:m2:the-fx-options-market:barrier}{options à barrière}.
\end{dated}

\section{Tutoriel~: des cotations au smile}\label{tut:m2:the-fx-options-market:smile}

\begin{tutorial}
\textbf{Objectif.} Pricer les options de change, calculer leurs deltas dans les quatre conventions,
retrouver les strikes à 25 delta et à la monnaie à partir des cotations des courtiers, construire le smile,
et calculer la couverture que libère une \omterm{def:m2:the-fx-options-market:barrier}{option knock-out} à sa barrière. \textbf{État final~:}
les \cref{fig:m2:the-fx-options-market:smile,fig:m2:the-fx-options-market:barrier},
les \cref{ex:m2:the-fx-options-market:strikes,ex:m2:the-fx-options-market:jpy} et les
nombres du problème du week-end.
\begin{tutsteps}
\item \textbf{Deltas et strikes}~: les quatre conventions de delta et le strike
associé à un delta donné.
\omcode{code/firm/fxsmile/firm_fxsmile.py}{32}{59}{Deltas spot, forward et ajustés de la prime, et strike associé à un delta donné.}\label{lst:m2:the-fx-options-market:delta}
\item \textbf{À la monnaie et le smile}~: le \omterm{def:m2:the-fx-options-market:dns}{straddle delta neutre}, les
volatilités simplifiées des ailes et un smile quadratique.
\omcode{code/firm/fxsmile/firm_fxsmile.py}{62}{81}{Strike à la monnaie, volatilités des ailes et smile à trois points.}\label{lst:m2:the-fx-options-market:smile}
\item \textbf{Exécuter} \texttt{fxopt\_demo.smile} pour chaque convention,
\texttt{fxopt\_demo.barrier()} et \texttt{fig\_fxopt.py}.
\end{tutsteps}
\textbf{Pour aller plus loin.} Remplacer la formule simplifiée par une calibration qui
reproduit exactement le \omterm{def:m2:the-fx-options-market:rr}{strangle} de marché, et comparer les strikes des ailes~; puis valoriser la
\omterm{def:m2:the-fx-options-market:barrier}{option knock-out} avec la volatilité du smile à la barrière plutôt qu'une volatilité plate.
\end{tutorial}

\section{Construction~: le convertisseur de cotations en smile}\label{bld:m2:the-fx-options-market:fxsmile}

\begin{build}
\bfield{Rôle} Le desk d'options de change de la firme miniature reçoit des cotations de courtiers en
ATM, RR et BF, et doit pricer des strikes isolés et des barrières de façon cohérente avec eux
et les couvrir avec la quantité de spot appropriée.
\bfield{Interface} \texttt{gk(s, k, t, rd, rf, vol, phi)}~; \texttt{delta(\ldots,
kind)} où kind prend l'une des valeurs spot, forward, spot\_pa, forward\_pa~;
\texttt{strike\_from\_delta}~; \texttt{atm\_dns}~; \texttt{smile\_vols(atm, rr, bf)}~;
\texttt{quadratic\_smile(points)}~; \texttt{down\_and\_out\_call}~;
\texttt{barrier\_delta}.
\bfield{Règles} Taux à capitalisation continue~; le prix d'exercice du call ajusté de la prime
doit être fixé au-dessus du maximum du delta~; la formule simplifiée pour les ailes, signalée comme
une approximation.
\bfield{Tests d'acceptation} \texttt{code/firm/fxsmile/tests/}~: parité call-put~;
allers-retours strike-delta dans les quatre conventions pour les calls et les puts~;
l'exemple publié d'ajustement de la prime~; strikes delta-neutres~; les points d'ancrage
du smile~; la valeur de l'option à la barrière et loin d'elle.
\bfield{Pour aller plus loin} Calibration sur le \omterm{def:m2:the-fx-options-market:rr}{strangle} de marché~; interpolation du smile en delta
(SABR ou vanna--volga, le livre~5)~; valorisation des barrières sous le smile.
\end{build}

\begin{omsources}
\item D.~Reiswich et U.~Wystup, «~FX volatility smile construction~», Centre for
Practical Quantitative Finance, document de travail 20, Frankfurt School of Finance and
Management, 2010 (conventions d'après I.~Clark, \emph{Foreign Exchange Option
Pricing}).
\item Bank for International Settlements, Triennial Central Bank Survey 2025.
\end{omsources}

\section{Exercices}

\begin{exercise}[$\star$]\label{exo:m2:the-fx-options-market:1}
Le \omterm{def:m2:the-fx-options-market:rr}{risk reversal} EURUSD à 25 delta vaut $-0.5$ et le \omterm{def:m2:the-fx-options-market:rr}{butterfly} 0.2, ATM 7~\%. Donner les
volatilités des ailes par la formule simplifiée, et dire quelle option est la plus chère.
\end{exercise}

\begin{exercise}[$\star$]\label{exo:m2:the-fx-options-market:2}
Indiquer quelle convention de delta le marché utilise pour EURUSD, USDJPY et USDBRL, et
pourquoi la devise de la prime importe.
\end{exercise}

\begin{exercise}[$\star$]\label{exo:m2:the-fx-options-market:3}
Un call en euros a un \omterm{def:m2:the-fx-options-market:delta}{delta spot} de 60~\% et coûte 73\,669 euros pour
1~million d'euros, payés en euros. Calculer combien d'euros le vendeur achète pour se couvrir.
\end{exercise}

\begin{exercise}[$\star\star$]\label{exo:m2:the-fx-options-market:4}
Donner les strikes à 25 delta du call et du put EURUSD à trois mois, avec les \omterm{def:m2:the-fx-options-market:delta}{deltas spot} et
forward usuels, en utilisant les cotations de l'\cref{exo:m2:the-fx-options-market:1}.
\end{exercise}

\begin{exercise}[$\star\star$]\label{exo:m2:the-fx-options-market:5}
Pourquoi le strike delta neutre ajusté de la prime se situe-t-il sous le forward, et le
strike usuel se situe-t-il au-dessus~?
\end{exercise}

\begin{exercise}[$\star\star$]\label{exo:m2:the-fx-options-market:6}
Pourquoi un call down-and-out coûte-t-il moins cher que le call vanille, et quand la
différence est-elle importante~?
\end{exercise}

\begin{exercise}[$\star\star\star$]\label{exo:m2:the-fx-options-market:7}
\emph{Codage.} Avec \texttt{strike\_from\_delta}, donner les strikes USDJPY à 25 delta
avec les \omterm{def:m2:the-fx-options-market:delta}{deltas forward}, et les comparer aux strikes ajustés spot et ajustés de la prime de
l'\cref{ex:m2:the-fx-options-market:jpy}.
\end{exercise}

\begin{exercise}[$\star\star\star$]\label{exo:m2:the-fx-options-market:8}
\emph{Repérer l'erreur.} «~Le \omterm{def:m2:the-fx-options-market:rr}{risk reversal} est négatif, donc le marché anticipe une baisse
du dollar face au yen.~» La corriger.
\end{exercise}

\section{Problème~: La défense de la barrière}

\begin{problem}[{Problème du week-end --- ce qu'un knock-out envoie sur le marché au comptant}]\label{pb:m2:the-fx-options-market:1}
Une entreprise achète à un dealer un call EURUSD down-and-out à trois mois sur
500~millions d'euros, strike 1.1464, barrière 1.1200, spot 1.1464, avec les taux de
l'\cref{ex:m2:the-fx-options-market:strikes} et une volatilité de 7~\%. Le dealer
couvre son delta.

\medskip\noindent\textbf{Partie I --- L'option.}
\begin{enumerate}
\item Donner la valeur du call vanille, en pips de notionnel.
\item Donner la valeur du knock-out et l'économie réalisée par l'entreprise.
\item Donner le delta actuel du knock-out, et la couverture du dealer.
\item Pourquoi le delta du knock-out est-il plus grand que celui du vanille~?
\item Donner le delta juste au-dessus de la barrière, et celui du vanille à cet endroit.
\end{enumerate}

\medskip\noindent\textbf{Partie II --- À la barrière.}
\begin{enumerate}[resume]
\item Que détient le dealer juste avant que le spot ne touche 1.1200~?
\item Que doit-il échanger au moment du contact, et dans quel sens~?
\item Comment cette transaction affecte-t-elle le spot~?
\item Que ferait un dealer acheteur d'une telle barrière lorsque le spot s'en approche~?
\item Pourquoi d'autres intervenants pourraient-ils opérer vers une grande barrière connue~?
\end{enumerate}

\medskip\noindent\textbf{Partie III --- La gestion.}
\begin{enumerate}[resume]
\item Comment le dealer peut-il étaler le dénouement au lieu de le réaliser au contact~?
\item Quel en est le coût, et quel risque subsiste-t-il~?
\item Pourquoi le gamma du dealer près de la barrière rend-il la couverture quotidienne coûteuse~?
\item Comment un smile de volatilité modifierait-il la valeur du knock-out~?
\item Que doit comprendre l'entreprise de l'option qu'elle a achetée~?
\end{enumerate}

\medskip\noindent\textbf{Partie IV --- Jugement.}
\begin{enumerate}[resume]
\item Défendre une barrière constitue-t-il une manipulation de marché~?
\item Pourquoi les barrières se regroupent-elles autour de nombres ronds~?
\item Quelles informations sur les barrières un dealer doit-il garder confidentielles~?
\item Énoncer le \emph{résultat nommé}~: le spot auquel le dealer doit vendre à la barrière.
\item En une phrase~: pourquoi les \omterm{def:m2:the-fx-options-market:barrier}{options à barrière} font-elles bouger le spot~?
\end{enumerate}
\end{problem}

\section{Questions d'entretien}

\begin{interviewq}[$\star$ \iqroles{trader, chercheur}]\label{iq:m2:the-fx-options-market:1}
Que vous apprennent la volatilité à la monnaie, le \omterm{def:m2:the-fx-options-market:rr}{risk reversal} à 25 delta et le \omterm{def:m2:the-fx-options-market:rr}{butterfly}
sur le smile d'une paire de devises~?
\end{interviewq}

\begin{interviewq}[$\star$ \iqroles{développeur, trader}]\label{iq:m2:the-fx-options-market:2}
Qu'est-ce qu'un \omterm{def:m2:the-fx-options-market:delta}{delta ajusté de la prime}, et dans quels cas l'utilise-t-on~?
\end{interviewq}

\begin{interviewq}[$\star\star$ \iqroles{chercheur}]\label{iq:m2:the-fx-options-market:3}
Étant donnés les cotations ATM, RR et BF, comment trouvez-vous le strike du put à 25 delta~?
\end{interviewq}

\begin{interviewq}[$\star\star$ \iqroles{trader}]\label{iq:m2:the-fx-options-market:4}
Vous êtes vendeur d'un gros call down-and-out et le spot baisse vers la barrière.
Que faites-vous~?
\end{interviewq}

\begin{interviewq}[$\star\star$ \iqroles{chercheur, trader}]\label{iq:m2:the-fx-options-market:5}
Pourquoi définit-on le niveau à la monnaie en FX comme le \omterm{def:m2:the-fx-options-market:dns}{straddle delta neutre} plutôt que
comme le forward~?
\end{interviewq}

\begin{interviewq}[$\star\star\star$ \iqroles{développeur}]\label{iq:m2:the-fx-options-market:6}
Concevez un service de surface de volatilité pour cinquante paires de devises, qui reçoit les
cotations des courtiers et sert les volatilités par strike et par échéance aux moteurs de pricing.
\end{interviewq}
